\normalsize
\section{Background}
\label{sec:background}

We first provide an introduction to LoRA fine-tuning (§\ref{subsec:background-llm-lora-finetuning}), and then present an overview of current system-level optimizations for fine-tuning (§\ref{subsec:background-system-optimization}).

\begin{table}[t]
    \centering
    \caption{Notation for the LoRA fine-tuning in this paper.}
     \begin{tabularx}{\linewidth}{c X}
     \toprule
     Symbol & Description \\
     \midrule
     $r$ & LoRA rank \\
     $m$ & Number of tokens (batch size $\times$ seq length) \\
     $k, n$ & Input/Output dimension of the weight matrix \\
    %  $n$ & Output dimension of the weight matrix \\
     \midrule
     $W$ & Base model weights. Size: $(k, n)$. \\
     $A$ & First LoRA weights. Size: $(k, r)$ \\
     $B$ & Second LoRA weights. Size: $(r, n)$ \\
     $X$ & Input tensor. Size: $(m, k)$ \\
     $\widehat{X}$ & Input tensor after dropout. Size: $(m, k)$ \\
     $Y$ & Output tensor. Size: $(m, n)$ \\
     \bottomrule
     \end{tabularx}
    \label{tab:background-tab-1-notation}
\end{table}

\subsection{LLM LoRA Fine-tuning}
\label{subsec:background-llm-lora-finetuning}

Training an LLM from scratch demands substantial amounts of data, millions of GPU hours, and significant costs~\cite{llama3.1-training-cost}.
Fine-tuning pre-trained LLMs, such as LLaMa~\cite{llama3,llama3.1} and Qwen~\cite{bai2023qwen,yang2024qwen2-5} is thus more practical.
Fine-tuning preserves pre-trained capabilities while adapting the model to specialized downstream tasks~\cite{devlin2018bert}.
While fine-tuning reduces the data needs and training iterations, traditional full-model fine-tuning still requires similarly large GPU memory, due to the vast number of trainable parameters.

PEFT methods~\cite{hu2022lora,li2021prefix-tuning,lester2021prompt-tuning} address the memory requirements by freezing pre-trained parameters (keeping them non-updatable) and only training a small set of newly introduced parameters called \textit{adapters}.
Among these methods, LoRA~\cite{hu2022lora} is currently most widely used due to its simplicity and effectiveness. As illustrated in Figure~\ref{fig:background-lora-illustration}, LoRA injects two small trainable matrices alongside each pretrained weight matrix.
%
Formally, for a pretrained weight matrix $W \in \mathbb{R}^{k \times n}$, LoRA introduces two low-rank matrices $A \in \mathbb{R}^{k \times r}$ and $B \in \mathbb{R}^{r \times n}$, where LoRA rank $r \ll k,n$, combined as:
\begin{equation}
Y = X W + \alpha S B = X W + \alpha (\widehat{X} A) B
\end{equation}
where $\widehat{X} \in \mathbb{R}^{m \times k}$ is the input after dropout of $X \in \mathbb{R}^{m \times k}$, $S=\widehat{X} A$ is the intermediate result, $m$ is the number of aggregated tokens (batch size multiplied by sequence length), and $\alpha$ is a constant scalar for scaling.
Table~\ref{tab:background-tab-1-notation} summarizes key notations.

When fine-tuning LLMs, LoRA is applied to linear layers, replacing each original layer of dimensions $k \times n$ with a LoRA-equipped version. 
Assuming half-precision training with full-precision optimizer~\cite{megatron-1,zero}, memory usage of model states per linear layer decreases from $16nk$ bytes to $2nk + 32r(n+k)$ bytes. 
Since $r$ is much smaller than $n$ and $k$, the memory footprint of trainable parameters, gradients, and optimizer states is negligible compared to the pre-trained weights.
For instance, with $n$=$k$=4096 and $r$=16, LoRA parameters account for just about 0.39\% of the pre-trained weights. 
This dramatically reduces the memory required for gradients and optimizer states, decreasing memory demands by nearly $8\times$ compared to full-model fine-tuning.

\subsection{System Optimization for Fine-tuning}
\label{subsec:background-system-optimization}

Despite algorithmic advancements in LoRA~\cite{dettmers2023qlora,kopiczko2023vera,liu2024dora}, and optimizations for LoRA serving~\cite{chen2024punica,sheng2023s-lora,wu2024dlora}, existing system optimizations for LoRA fine-tuning primarily reuse techniques originally developed for LLM pre-training, including parallelization~\cite{pytorch-distributed,tensorflow,zero,pytorch-fsdp,megatron-1,zheng2022alpa}, kernel fusion~\cite{tvm,hidet,ansel2024pytorch-compile,dao2022flashattention,hsu2024liger,tillet2019triton,ashari2015optimizing}, and data packing~\cite{kundu2024hf-packing-with-flash-attention,wang2024packing}.

\myparagraph{Parallelization}
Parallelization partitions the computation and memory usage of the training process across multiple GPUs.
Data parallelism (DP)~\cite{pytorch-distributed,tensorflow} replicates the model across GPUs and partitions data batches. 
Fully Sharded Data Parallelism (FSDP or ZeRO-3)~\cite{zero,pytorch-fsdp} partitions model states and communicates them only when necessary, significantly reducing memory usage.
Tensor Parallelism (TP)~\cite{megatron-1, megatron-2} splits linear layers across GPUs and merges partial results via communication. 
Pipeline Parallelism (PP)~\cite{gpipe,pipedream,dapple} divides the model into sequential stages executed in a pipeline manner, reducing communication overhead but potentially introducing idle time (pipeline bubbles).

\myparagraph{Kernel Fusion}
Kernel fusion improves efficiency by merging multiple operations into fewer GPU kernels, reducing memory transfers and kernel launch overhead~\cite{tvm,hidet,ansel2024pytorch-compile,dao2022flashattention,hsu2024liger,tillet2019triton,ashari2015optimizing}. 
Recent techniques like Flash-Attention~\cite{dao2022flashattention} and element-wise kernel fusion~\cite{tillet2019triton,hsu2024liger} significantly improve performance by fusing frequently used operations in LLMs.
For LoRA, specialized kernels~\cite{chen2024punica,sheng2023s-lora} for multi-adapter serving are proposed, but these kernels do not directly improve fine-tuning throughput, as we discuss in Section~\ref{sec:motivation}.

\begin{figure}[t]
    \centering
    \includegraphics[width=0.995\linewidth]{figures/background-fig-2-padding-and-packing.pdf}
    \caption{Comparison of (a) traditional batch padding, (b) dataset pre-packing, and (c) batch on-the-fly packing methods for LoRA fine-tuning of LLMs.}
    \Description{Padding and packing for LoRA fine-tuning.}
    \label{fig:background-padding-and-packing}
\end{figure}

\myparagraph{On-the-fly Data Packing}
When fine-tuning LLMs, training data often consists of token sequences of variable lengths. 
As depicted in Figure~\ref{fig:background-padding-and-packing}(a), traditional padding aligns shorter samples with padding tokens, causing wasted computations. 
Figure~\ref{fig:background-padding-and-packing}(b) shows the dataset pre-packing, which forms fixed-length batches in advance, but introduces variable sample counts per batch, potentially affecting training stability and randomness if not properly handled~\cite{kundu2024hf-packing-with-flash-attention}.
In contrast, on-the-fly packing (Figure~\ref{fig:background-padding-and-packing}(c)) dynamically concatenates samples within each batch, avoiding wasted computations while maintaining deterministic training samples per batch. Given its effectiveness and wide usage, we adopt on-the-fly packing throughout our work.
